\documentclass[11pt]{article}

\usepackage[letterpaper,margin=1in]{geometry}
\usepackage[T1]{fontenc}
\usepackage[utf8]{inputenc}
\usepackage{lmodern}
\usepackage{microtype}
\usepackage{amsmath,amssymb,mathtools,bm}
\usepackage{enumitem}
\usepackage{xcolor}
\usepackage{booktabs}
\usepackage{array}
\usepackage{hyperref}
\usepackage[most]{tcolorbox}
\usepackage{fancyhdr}
\usepackage{titlesec}
\usepackage{parskip}
\usepackage{needspace}

\definecolor{qftblue}{HTML}{245A8D}
\definecolor{qftgreen}{HTML}{2F6B4F}
\definecolor{qftred}{HTML}{A23B3B}
\definecolor{qftgray}{HTML}{F4F6F8}

\hypersetup{
  colorlinks=true,
  linkcolor=qftblue,
  urlcolor=qftblue,
  citecolor=qftgreen,
  pdftitle={PHYS 662 Lecture 12: Weyl Spinors and Invariant Tensors},
  pdfauthor={Nima Lashkari}
}

\pagestyle{fancy}
\fancyhf{}
\lhead{PHYS 662: Quantum Field Theory I}
\rhead{Lecture 12}
\cfoot{\thepage}
\setlength{\headheight}{14pt}

\titleformat{\section}{\Large\bfseries\color{qftblue}}{\thesection.}{0.6em}{}
\titleformat{\subsection}{\large\bfseries\color{qftgreen}}{\thesubsection}{0.6em}{}
\titleformat{\subsubsection}{\normalsize\bfseries\color{qftred}}{\thesubsubsection}{0.6em}{}

\setlist[itemize]{leftmargin=2em,itemsep=0.25em,topsep=0.4em}
\setlist[enumerate]{leftmargin=2em,itemsep=0.25em,topsep=0.4em}

\newcommand{\dd}{\mathrm{d}}
\newcommand{\R}{\mathbb{R}}
\newcommand{\Lag}{\mathcal{L}}
\newcommand{\cM}{\mathcal{M}}
\newcommand{\SO}{\mathrm{SO}}
\newcommand{\OO}{\mathrm{O}}

\newtcolorbox{keyidea}[1][]{
  enhanced, breakable, colback=qftgray, colframe=qftblue,
  boxrule=0.8pt, arc=1.5mm, title={#1}, fonttitle=\bfseries,
  coltitle=white
}
\newtcolorbox{remarkbox}[1][]{
  enhanced, breakable, colback=white, colframe=qftgreen,
  boxrule=0.7pt, arc=1.5mm, title={#1}, fonttitle=\bfseries,
  coltitle=white
}
\newtcolorbox{warningbox}[1][]{
  enhanced, breakable, colback=white, colframe=qftred,
  boxrule=0.7pt, arc=1.5mm, title={#1}, fonttitle=\bfseries,
  coltitle=white
}


\newcommand{\Th}{\Theta}
\newcommand{\Eul}{\mathcal E}
\newcommand{\Lie}{\mathscr L}

\newcommand{\C}{\mathbb C}
\newcommand{\Z}{\mathbb Z}
\newcommand{\SU}{\mathrm{SU}}
\newcommand{\UU}{\mathrm{U}}
\newcommand{\GL}{\mathrm{GL}}
\newcommand{\SL}{\mathrm{SL}}
\newcommand{\Sp}{\mathrm{Sp}}
\newcommand{\tr}{\operatorname{tr}}
\newcommand{\Ad}{\operatorname{Ad}}
\newcommand{\ad}{\operatorname{ad}}
\newcommand{\lieg}{\mathfrak g}
\newcommand{\lieh}{\mathfrak h}
\newcommand{\su}{\mathfrak{su}}
\newcommand{\so}{\mathfrak{so}}
\newcommand{\slie}{\mathfrak{sl}}


\newcommand{\cH}{\mathcal H}
\newcommand{\Spin}{\operatorname{Spin}}
\newcommand{\id}{\mathrm{id}}
\newcommand{\Lor}{\SO^+(1,3)}
\newcommand{\eps}{\varepsilon}
\newcommand{\one}{\mathbf 1}

\newcommand{\smu}{s}
\newcommand{\tmu}{\widetilde s}
\newcommand{\Eps}{E}
\newcommand{\pslash}{\slashed p}
\newcommand{\pL}{P_L}
\newcommand{\pR}{P_R}
\newcommand{\hc}{\mathrm{h.c.}}
\newcommand{\diag}{\operatorname{diag}}
\newcommand{\vol}{\mathrm{vol}}
\newcommand{\supp}{\operatorname{supp}}
\begin{document}
\pagenumbering{gobble}
\begin{titlepage}
\centering
\vspace*{1.2in}
{\scshape\Large PHYS 662 --- Quantum Field Theory I\par}
\vspace{0.35in}
{\Huge\bfseries\color{qftblue} Lecture 12\par}
\vspace{0.15in}
{\LARGE\bfseries Weyl Spinors\par}
\vspace{0.06in}
{\LARGE\bfseries and Invariant Tensors\par}
\vspace{0.16in}
{\large\color{qftgreen} From Two Spin Labels to Covariant Bilinears\par}
\vspace{0.6in}
{\large Nima Lashkari\par}
{\large Fall 2026\par}
\vfill
\begin{minipage}{0.86\textwidth}\small
The two Lorentz spin labels lead to two inequivalent Weyl representations.
We make their transformations explicit, distinguish conjugation from
duality, and construct vectors and invariant tensors from spinors.
The resulting calculus of scalar and vector bilinears, including
Grassmann signs, supplies the ingredients for the actions of Lecture 13.

\medskip
Course page: \url{https://nimalashkari.com/quantum-field-theory-2026-part-i/}
\end{minipage}
\vspace{0.5in}
\end{titlepage}
\pagenumbering{roman}
\begingroup\small\tableofcontents\endgroup
\newpage
\pagenumbering{arabic}

\section*{Roadmap}
\addcontentsline{toc}{section}{Roadmap}
Lecture 11 classified finite-dimensional complex Lorentz representations
by \((j_L,j_R)\). We now make the smallest nontrivial representations
concrete and learn to combine them into quantities with definite
Lorentz transformation laws.
\begin{enumerate}
\item Construct the two Weyl actions and explain how parity exchanges them.
\item Distinguish conjugation, duality, and raising or lowering a spinor index.
\item Build scalars, vectors, and tensor representations from spinors.
\item Derive sigma-matrix covariance and the Grassmann signs needed for actions.
\item Use covariance to write a first-order equation and solve a massless plane wave.
\end{enumerate}
This lecture develops the algebraic ingredients. Lecture 13 will derive
the Weyl, Dirac, and Majorana actions from them; Lecture 14 will connect
the covariant fields to one-particle states.

\begin{keyidea}[Conventions continued from Lecture 11]
We use \(\eta_{\mu\nu}=\diag(1,-1,-1,-1)\), \(\hbar=c=1\), and
\(\eps_{123}=1\). The fundamental spinor transforms as
\[
 \xi'(x)=A\xi(\Lambda^{-1}x),\qquad
 A=\exp\!\left[-\frac{i}{2}\bm\theta\cdot\bm\sigma
                    -\frac12\bm\rho\cdot\bm\sigma\right].
\]
It has \(J_i=\sigma_i/2\), \(K_i=-i\sigma_i/2\), and is called
\((\tfrac12,0)\) in our convention. Its conjugate dual transforms with
\(\bar A=(A^\dagger)^{-1}\) and is \((0,\tfrac12)\).
The active boost along \(z\) has
\(t'=t\cosh\rho-z\sinh\rho\).
We denote the two Pauli four-tuples by
\[
 s^\mu=(\one,\bm\sigma),\qquad
 \widetilde s^\mu=(\one,-\bm\sigma).
\]
With this choice the kinetic matrix for \(\xi\) is \(s^\mu\).
Some texts exchange the two chiral names and use \(\bar\sigma^\mu\)
for the corresponding kinetic matrix. Explicit matrices, rather than
the words left and right, fix every sign below.
\end{keyidea}

\section{The two Weyl representations}
\subsection{Explicit matrices and an elementary boost}
In \((\tfrac12,0)\), take \(L_i=\sigma_i/2\) and \(R_i=0\).
In \((0,\tfrac12)\), interchange the roles. The relations
\(J_i=L_i+R_i\), \(K_i=-i(L_i-R_i)\) from Lecture 11 give
\begin{equation}
 \begin{array}{c|cc}
 &J_i&K_i\\ \hline
 (\tfrac12,0)&\sigma_i/2&-i\sigma_i/2\\[2pt]
 (0,\tfrac12)&\sigma_i/2&+i\sigma_i/2
 \end{array}
 \label{eq:weyl-generators}
\end{equation}
The rotations agree and the boosts have opposite signs. For the same
covering-group element \(A\), the two actions are
\begin{equation}
 \boxed{\psi_L\mapsto A\psi_L,\qquad
 \psi_R\mapsto(A^\dagger)^{-1}\psi_R.}
 \label{eq:weyl-actions}
\end{equation}
The second action is a representation because taking the adjoint and
then the inverse reverses the product order twice. For a pure rotation
\(A\) is unitary, so \((A^\dagger)^{-1}=A\). For the boost \(A_z(\rho)=e^{-\rho\sigma_3/2}\),
\[
 \begin{pmatrix}\psi_{L1}\\\psi_{L2}\end{pmatrix}
 \mapsto
 \begin{pmatrix}e^{-\rho/2}\psi_{L1}\\e^{\rho/2}\psi_{L2}\end{pmatrix},
 \qquad
 \begin{pmatrix}\psi_{R1}\\\psi_{R2}\end{pmatrix}
 \mapsto
 \begin{pmatrix}e^{\rho/2}\psi_{R1}\\e^{-\rho/2}\psi_{R2}\end{pmatrix}.
\]
This both displays their inequivalence and makes nonunitarity visible:
the ordinary sum of squared component magnitudes changes under a boost.

\subsection{Why finite-dimensional boosts are nonunitary}
Why are the nontrivial finite-dimensional Lorentz representations
nonunitary? Suppose all \(J_i,K_i\) were finite Hermitian matrices.
Cyclicity of the trace and the rotation--boost algebra of Lecture 11 imply
\[
 -i\tr J_3^2=\tr([K_1,K_2]J_3)
 =\tr(K_1[K_2,J_3])=i\tr K_1^2.
\]
Both traces of squares are nonnegative, so both vanish. Cyclic
permutations give \(J_i=K_i=0\) for every \(i\).
\begin{keyidea}[Which representations have we classified?]
Every nontrivial finite-dimensional complex representation of the
connected Lorentz group or its cover is nonunitary. This is compatible
with unitary quantum symmetry because a physical particle's momentum
states form an infinite-dimensional space. The labels \((j_L,j_R)\)
classify covariant field components; they do not by themselves specify
a particle's mass, physical polarizations, or equations of motion.
\end{keyidea}

\subsection{Parity exchanges the two labels}
Under spatial parity, angular momentum is unchanged but boost charge
changes sign. The boost charge of Lecture 11,
\(K_i=\int\dd^3x\,x^iT^{00}-tP^i\), makes this explicit:
\[
 P:\quad J_i\mapsto J_i,\qquad K_i\mapsto-K_i.
\]
Hence \(L_i\leftrightarrow R_i\), and parity maps \((j_L,j_R)\) to
\((j_R,j_L)\). A single left Weyl representation cannot support parity
as a linear operation within its two-component space. The direct sum
\((\tfrac12,0)\oplus(0,\tfrac12)\) can. For
\(\Psi=(\psi_L,\psi_R)^T\), a possible component matrix is
\[
 \mathsf P=\begin{pmatrix}0&I_2\\I_2&0\end{pmatrix},\qquad
 \mathsf P\bm J\mathsf P^{-1}=\bm J,\qquad
 \mathsf P\bm K\mathsf P^{-1}=-\bm K.
\]
The field transformation also evaluates the old field at \((t,-\bm x)\),
and can include an intrinsic-parity phase. Whether this operation is
a symmetry of the Lagrangian depends on its couplings; representation
theory alone does not require parity invariance. Chirality labels the
Lorentz representation. Helicity is the spin projection along momentum,
and relating the two requires the field equations and a particle-state analysis.


\Needspace{8\baselineskip}
\section{Conjugation, duality, and the invariant epsilon tensor}
\subsection{Four related transformation laws}
Let \(\xi_\alpha\) and \(\chi_\alpha\) be columns vector with two components, with
\(\alpha=1,2\). Suppressing the coordinate argument for a moment,
\begin{equation}
 \xi\mapsto A\xi,\qquad
 \xi^*\mapsto A^*\xi^*,\qquad
 \xi^\dagger\mapsto\xi^\dagger A^\dagger.
 \label{eq:conjugates}
\end{equation}
A {\it dual} vector transforms instead with \(A^{-T}=(A^T)^{-1}\), so that its
transpose can contract with \(\xi\). A {\it conjugate dual} vector 
transforms with \((A^\dagger)^{-1}\). Complex conjugation and taking
the dual are different operations; in particular \(A^*\) is not
generally \((A^\dagger)^{-T}\) unless the matrices $A$ are unitary.

Introduce the antisymmetric matrix
\begin{equation}
 E=(\eps_{\alpha\beta})=
 \begin{pmatrix}0&1\\-1&0\end{pmatrix},\qquad
 E^{-1}=-E,\qquad E^2=-\one.
 \label{eq:E}
\end{equation}
For any two-by-two matrix,
\(A^T E A=(\det A)E\). Hence for \(A\in\SL(2,\C)\),
\begin{equation}
 A^T E A=E,\qquad E A=A^{-T}E,\qquad
 E A^*=(A^\dagger)^{-1}E=\bar{A}E.
 \label{eq:intertwiner}
\end{equation}
The first identity can be checked by multiplying a matrix
\(\begin{psmallmatrix}a&b\\c&d\end{psmallmatrix}\) by \(E\):
the only surviving coefficient is \(ad-bc\). Thus the invariant
epsilon tensor identifies a representation with its dual.
It also turns a complex-conjugate spinor into a conjugate-dual spinor. In mathematical language, matrices $E$ are intertwiners that intertwine (identify) the two representations.



\subsection{Raising and lowering indices}
Set \(\eps^{\alpha\beta}=-\eps_{\alpha\beta}\), as matrices, and define
\begin{equation}
 \xi^\alpha=\eps^{\alpha\beta}\xi_\beta,\qquad
 \xi_\alpha=\eps_{\alpha\beta}\xi^\beta,
 \qquad \eps^{\alpha\gamma}\eps_{\gamma\beta}
       =\delta^\alpha{}_{\beta}.
 \label{eq:raise}
\end{equation}
In components, \(\xi^1=-\xi_2\) and \(\xi^2=\xi_1\).
A dotted index keeps track of a complex-conjugate representation:
\(\xi^*_{\dot\alpha}\) transforms with \(A^*\). Raising its dotted
index with the same numerical epsilon matrix gives
\begin{equation}
 r^{\dot\alpha}=\eps^{\dot\alpha\dot\beta}\xi^*_{\dot\beta},
 \qquad r=-E\xi^*,\qquad r\mapsto\bar A r.
 \label{eq:r}
\end{equation}
The dot is a representation label, not another spatial index. It cannot
be removed by merely relabeling \(\dot\alpha\) as \(\alpha\).

The basic scalar contraction is
\begin{equation}
 B(\xi,\chi):=\xi^\alpha\chi_\alpha
    =\xi^T E\chi=\xi_1\chi_2-\xi_2\chi_1.
 \label{eq:B}
\end{equation}
Equation \eqref{eq:intertwiner} proves its invariance directly.
Writing the contraction as a matrix product is often the fastest way
to diagnose an index-sign error.

\begin{warningbox}[An index identity is not yet a statistics identity]
The invariant tensor \(E\) is antisymmetric. Whether a contraction
changes sign when the \emph{fields} are exchanged also depends on whether
their components commute or anticommute. These are two separate signs.
\end{warningbox}


\subsection{Self-duality is not equivalence to the conjugate}
Equation \eqref{eq:intertwiner} says that the fundamental is equivalent
to its dual, while its complex conjugate is equivalent to the opposite
Weyl representation. These are different equivalences.
No invertible complex-linear matrix intertwines the two Weyl actions:
on rotations it would commute with every Pauli matrix, hence be a scalar
by Schur's lemma. Such a scalar cannot intertwine their opposite,
nonzero boost generators. Restricted to \(\SU(2)\), the fundamental
is equivalent to its conjugate; the boosts obstruct extending that
equivalence to \(\SL(2,\C)\).

\begin{warningbox}[Representations are conjugate; fields can be independent]
The column \(r=-E\xi^*\) constructs an opposite-chirality field from
a chosen \(\xi\). It does not require an independently introduced
right Weyl field to equal that column. A Dirac field has independent
left and right components. A Majorana condition imposes a relation
between them when compatible with the other symmetries.
\end{warningbox}

\subsection{Why a Hermitian product is different}
For two fundamental columns, \(\xi^\dagger\chi\) transforms into
\(\xi^\dagger A^\dagger A\chi\). A boost along \(z\) changes its value
from \(1\) to \(e^{-\rho}\) when \(\xi=\chi=(1,0)^T\).
In contrast, for \(\xi=(1,0)^T\) and \(\chi=(0,1)^T\),
\(\xi^T E\chi=1\) stays fixed: the two components acquire opposite
boost factors. Self-duality supplies an invariant antisymmetric
bilinear form, not a positive Hermitian norm.

If \(r\) transforms with \((A^\dagger)^{-1}\), then \(r^\dagger\xi\)
is a scalar because \(r^\dagger\mapsto r^\dagger A^{-1}\).
This is distinct from the Hermitian product of two fundamental columns.
It also explains why conjugation must be tracked before deciding whether
a mixed-chirality expression is a scalar or a vector.

\section{Tensor products: what can spinors build?}
\subsection{Two spinors of the same chirality}
Angular-momentum addition in each complexified Lorentz factor gives
\begin{equation}
 (\tfrac12,0)\otimes(\tfrac12,0)=(0,0)\oplus(1,0).
 \label{eq:samespin}
\end{equation}
For a general tensor \(T_{\alpha\beta}\), the decomposition is explicit:
\begin{equation}
 T_{\alpha\beta}=S_{\alpha\beta}+E_{\alpha\beta}b,
 \quad S_{\alpha\beta}=S_{\beta\alpha},\quad
 b=\tfrac12(T_{12}-T_{21})
   =-\tfrac12\eps^{\alpha\beta}T_{\alpha\beta}.
 \label{eq:decomp}
\end{equation}
The symmetric part has three components and transforms as \((1,0)\);
the antisymmetric part has one component and is a scalar.
Two conjugate-chirality spinors similarly give
\((0,\tfrac12)\otimes(0,\tfrac12)=(0,0)\oplus(0,1)\).

This representation-theory statement is about an unrestricted tensor
product. For the product of one Grassmann field with itself at the same
point, \(\xi_\alpha\xi_\beta\) is antisymmetric, so only the scalar
part remains. For two distinct fields, a symmetric spinor tensor need
not vanish.

\subsection{A vector as a bispinor}
Mixed chirality gives
\begin{equation}
 (\tfrac12,0)\otimes(0,\tfrac12)=(\tfrac12,\tfrac12).
 \label{eq:mixed}
\end{equation}
The same vector representation appears in the Hermitian matrix
\begin{equation}
 V=v^0\one+v^i\sigma_i=v_\mu\widetilde s^\mu,
 \qquad V'=A V A^\dagger.
 \label{eq:vmap}
\end{equation}
For example, \(\xi\chi^\dagger\) has this transformation law; it is a
complex vector in general and becomes Hermitian when \(\chi=\xi\).
The inverse map and invariant norm are
\begin{equation}
 v^\mu=\tfrac12\tr(s^\mu V),\qquad
 \det V=(v^0)^2-|\bm v|^2.
 \label{eq:inversev}
\end{equation}
There is a factor of \(1/2\) in the inverse map because
\(\tr\one=2\). Hermiticity selects real vectors; arbitrary complex
two-by-two matrices describe the complexified vector representation.

\begin{remarkbox}[A commuting spinor gives a future null vector]
Take an ordinary complex column \(z=(a,b)^T\) and set \(X=zz^\dagger\).
Using Eq.~\eqref{eq:inversev},
\[
 v^0=\tfrac12(|a|^2+|b|^2),\quad
 v^1=\Re(a^*b),\quad v^2=\Im(a^*b),\quad
 v^3=\tfrac12(|a|^2-|b|^2).
\]
The matrix has rank one for \(z\ne0\), so \(\det X=0\) and \(v^0>0\).
Thus it describes a future null vector. For \(z=(1,0)^T\),
\(v=(1,0,0,1)/2\); the boost \(A_z(\rho)=e^{-\rho\sigma_3/2}\) scales this
vector by \(e^{-\rho}\). This is a concrete example of a vector made
from two spinor factors. Positivity in this example uses ordinary
complex numbers, not Grassmann-valued components.
\end{remarkbox}

\subsection{Rank-two tensors and two-forms}
Taking the tensor product of two vectors gives
\begin{equation}
 (\tfrac12,\tfrac12)\otimes(\tfrac12,\tfrac12)
  =(0,0)\oplus(1,0)\oplus(0,1)\oplus(1,1).
 \label{eq:ranktwo}
\end{equation}
The dimensions \(1+3+3+9=16\) have a familiar spacetime meaning:
\begin{equation}
 T_{\mu\nu}=\frac14\eta_{\mu\nu}T^\rho{}_{\rho}
  +\left(T_{(\mu\nu)}-\frac14\eta_{\mu\nu}T^\rho{}_{\rho}\right)
  +T_{[\mu\nu]}.
 \label{eq:tensordecomp}
\end{equation}
These are the trace transforming in $(0,0)$, symmetric traceless tensor transforming in $(1,1)$, and antisymmetric
two-form transforming as  \((1,0)\oplus(0,1)\) (it is reducible over \(\C\)).

To see why complexification matters, fix an orientation and define
\((*F)_{\mu\nu}=\tfrac12\eps_{\mu\nu\rho\sigma}F^{\rho\sigma}\).
With Lorentzian signature, \(*^2=-1\) on two-forms. The combinations
\begin{equation}
 F_+=\tfrac12(F-i*F),\qquad F_-=\tfrac12(F+i*F),
 \qquad *F_\pm=\pm iF_\pm
 \label{eq:selfdual}
\end{equation}
are the two irreducible three-dimensional pieces. Which piece is called
\((1,0)\) depends on the orientation and chiral naming convention.
For a real two-form, \(F_-=(F_+)^*\), so these are not two independent
real tensors. A symmetric spinor \(S_{\alpha\beta}\) is another way
to describe one of the complex pieces.

Under spatial rotations alone,
\begin{equation}
 (j_L,j_R)\big|_{\SU(2)_{\rm rot}}
     =\bigoplus_{j=|j_L-j_R|}^{j_L+j_R}j.
 \label{eq:rotationcontent}
\end{equation}
Thus a four-vector restricts to \(0\oplus1\), not \(0\otimes1\):
its time component is a rotational scalar and its spatial components
form a triplet. This is a statement about field components, not yet a
count of physical particle polarizations.

\section{Sigma matrices and Lorentz-covariant bilinears}
\subsection{The identities we actually need}
The Pauli multiplication rule
\(\sigma_i\sigma_j=\delta_{ij}\one+i\eps_{ijk}\sigma_k\)
implies
\begin{align}
 s^\mu\widetilde s^\nu+s^\nu\widetilde s^\mu
    &=2\eta^{\mu\nu}\one,\nonumber\\
 \widetilde s^\mu s^\nu+\widetilde s^\nu s^\mu
    &=2\eta^{\mu\nu}\one,\label{eq:sigmaclifford}\\
 \tr(s^\mu\widetilde s^\nu)&=2\eta^{\mu\nu}.
 \label{eq:sigmatrace}
\end{align}
where we have defined
\begin{align}
    s^\mu=({\bf 1},{\bf \sigma}),\qquad \tilde{s}^\mu=({\bf 1},-{\bf \sigma})
\end{align}
transforming as
\begin{align}
    A^\dagger s^\mu A=\Lambda^\mu_{:\nu} s^\nu
\end{align}
and the conjugate dual
\begin{align}
    \bar{A}^\dagger \tilde{s}^\mu \bar{A}=\Lambda^\mu_{:\nu} \tilde{s}^\nu\ .
\end{align}

For instance, the time-space anticommutator is
\(-\sigma_i+\sigma_i=0\), while the space-space one is
\(-\{\sigma_i,\sigma_j\}=-2\delta_{ij}\one\).
Contracting the first identity with a symmetric product \(p_\mu p_\nu\)
gives an especially useful result:
\begin{equation}
 (p\cdot s)(p\cdot\widetilde s)=p^2\one,
 \quad p\cdot s=p^0\one-\bm p\cdot\bm\sigma,
 \quad p\cdot\widetilde s=p^0\one+\bm p\cdot\bm\sigma.
 \label{eq:sigmap}
\end{equation}
The Lorentz contraction uses \(p_\mu\), not four entries with all plus
signs. One may also verify the completeness identity in ordinary
matrix indices \(a,b,c,d\):
\begin{equation}
 \sum_{\mu,\nu}\eta_{\mu\nu}(s^\mu)_{ab}(s^\nu)_{cd}
       =2 E_{ac}E_{bd}.
 \label{eq:completeness}
\end{equation}
Its origin is
\(\sum_i(\sigma_i)_{ab}(\sigma_i)_{cd}
=2\delta_{ad}\delta_{bc}-\delta_{ab}\delta_{cd}\).
Spinor rearrangement identities follow from this formula together
with the appropriate Grassmann signs.

\subsection{A vector bilinear and its covariance}
The map \eqref{eq:vmap} and its inverse imply
\begin{equation}
 A^\dagger s^\mu A=\Lambda^\mu{}_{\nu}s^\nu,
 \qquad
 \bar A^\dagger\widetilde s^\mu\bar A
       =\Lambda^\mu{}_{\nu}\widetilde s^\nu.
 \label{eq:covsigma}
\end{equation}
For the first identity, take \(\tfrac12\tr(s^\mu A V A^\dagger)\)
and use cyclicity of the trace. Since it equals
\(\Lambda^\mu{}_{\nu}v^\nu\) for every \(V\), the matrix identity
follows. Consequently,
\begin{equation}
 j^\mu(\xi,\chi)=\xi^\dagger s^\mu\chi
 \quad\Longrightarrow\quad
 j'^\mu(x)=\Lambda^\mu{}_{\nu}j^\nu(\Lambda^{-1}x).
 \label{eq:j}
\end{equation}
Moreover \((j^\mu(\xi,\chi))^\dagger=j^\mu(\chi,\xi)\), because
the matrices \(s^\mu\) are Hermitian. In particular
\(\xi^\dagger s^\mu\xi\) is a Hermitian vector bilinear.

For comparison, the two chiral spin generators can be written
\begin{align}
 S_L^{\mu\nu}&=\frac{i}{4}
   (\widetilde s^\mu s^\nu-\widetilde s^\nu s^\mu),\nonumber\\
 S_R^{\mu\nu}&=\frac{i}{4}
   (s^\mu\widetilde s^\nu-s^\nu\widetilde s^\mu).
 \label{eq:twogenerators}
\end{align}
Thus \(\tfrac12\eps_{ijk}S_{L,R}^{jk}=\sigma_i/2\), but
\(S_L^{i0}=-i\sigma_i/2\) and \(S_R^{i0}=+i\sigma_i/2\).
The rotation generators agree and the boost generators differ in sign.
Equivalently the epsilon-lowered matrices \(E\sigma_i\) are symmetric;
this exhibits the link between the three spin matrices and symmetric
two-index spinors.

\subsection{Spin transformation versus coordinate transformation}
A field changes both its components and its argument:
\begin{equation}
 \Phi'(x)=D(\Lambda)\Phi(\Lambda^{-1}x).
 \label{eq:fieldlaw}
\end{equation}
For \(\Lambda=\one+\omega+O(\omega^2)\) and
\(D=\one+\delta D+O(\omega^2)\), expansion gives
\begin{equation}
 \delta\Phi(x)=\delta D\,\Phi(x)
            -\omega^\mu{}_{\nu}x^\nu\partial_\mu\Phi(x).
 \label{eq:orbitalspin}
\end{equation}
The first term acts on spin indices and the second is the orbital
action on spacetime dependence. At \(x=0\), only the first survives.
This explains why transformations of a field at the origin expose its
finite-dimensional Lorentz representation without displaying the
orbital generators.

\section{Fermions and Grassmann signs}
\subsection{Why a scalar made from one spinor need not vanish}
Classical fermionic fields are taken to be Grassmann odd. Thus
\(\xi_\alpha\chi_\beta=-\chi_\beta\xi_\alpha\), including when
the two fields are the same. In a component calculation,
\begin{align}
 B(\chi,\xi)&=\chi_1\xi_2-\chi_2\xi_1
             =\xi_1\chi_2-\xi_2\chi_1=B(\xi,\chi),
             \label{eq:Bsym}\\
 B(\xi,\xi)&=2\xi_1\xi_2.\label{eq:selfB}
\end{align}
The minus sign from exchanging fields cancels the minus sign from
the antisymmetric invariant tensor. For commuting spinors the result
would instead be antisymmetric, and the self-contraction would vanish.

Hermitian conjugation reverses order:
\begin{equation}
 (XY)^\dagger=Y^\dagger X^\dagger,\qquad
 B(\xi,\xi)^\dagger=2\xi_2^*\xi_1^*
                         =-2\xi_1^*\xi_2^*.
 \label{eq:daggerB}
\end{equation}
No extra minus sign is inserted into the definition of conjugation.
A minus sign appears only if the resulting odd variables are then
reordered. Treat \(\xi\) and \(\xi^\dagger\) as independent variables
when varying a classical action.

\begin{remarkbox}[Statistics and the scope of this argument]
Lorentz covariance by itself does not choose commuting or
anticommuting fields. Ordinary wavefunctions and polarization spinors
have commuting complex components; products of quantum fields at the
same point require an ordering or regularization prescription. Grassmann variables are the classical language
that leads to fermionic canonical anticommutation relations. The
spin-statistics connection requires the further physical assumptions
of relativistic quantum field theory; it is not proved just by an
epsilon-tensor identity.
\end{remarkbox}


\section{A first-order equation from covariance}
The sigma matrices also identify a Lorentz-covariant differential
operator. The identity \(A^\dagger s^\mu A=\Lambda^\mu{}_{\nu}s^\nu\)
implies that \(s^\mu\partial_\mu\xi\) transforms with
\((A^\dagger)^{-1}\). To see this, note that $\sigma^\mu$ are just fixed matrices that do not transform, so
\begin{align}
    \sigma^\mu(\partial_\mu\xi)\to \sigma^\mu(\Lambda^{-1})^\mu_{\:\nu}\partial_\nu(A\xi),
\end{align}

Thus, setting this column to zero is a
Lorentz-invariant equation:
\[
 i s^\mu\partial_\mu\xi=0.
\]
This is the massless Weyl equation. Lecture 13 will derive it from an
action and show how a mass modifies it. Here it provides a first
application of the spinor calculus.

Take \(\xi(x)=u\,e^{-ip\cdot x}\) with \(p^0>0\), where \(u\)
is an ordinary complex coefficient column. Then
\begin{equation}
 (p^0-\bm\sigma\cdot\bm p)u=0.
 \label{eq:masslessplane}
\end{equation}
For nonzero \(u\), the determinant condition gives \(p^2=0\), and
\(\bm\sigma\cdot\widehat{\bm p}\,u=u\). This is one spin
polarization at fixed positive-energy momentum. In the explicit
convention of these lectures its helicity is \(+1/2\); the opposite
chiral equation has helicity \(-1/2\). The words left and right are
convention dependent here, but the eigenvalue of
\(\bm J\cdot\widehat{\bm p}\) is not. We will distinguish these
wavefunctions from particle and antiparticle states in Lecture 14.


\Needspace{12\baselineskip}
\section*{Summary}
\addcontentsline{toc}{section}{Summary}
The Weyl representations have the same rotation matrices and opposite
boost generators. Parity exchanges them, while complex conjugation,
duality, and epsilon-index operations have distinct meanings.
The invariant epsilon tensor builds scalar contractions; mixed
chirality builds vectors; symmetric equal-chirality tensors describe
the complex two-form pieces. Sigma-matrix identities make these
transformation laws explicit.

For fermionic fields, the Grassmann exchange sign combines with the
antisymmetry of epsilon so that \(B(\xi,\chi)=B(\chi,\xi)\) and
\(B(\xi,\xi)=2\xi_1\xi_2\) need not vanish. These are precisely the
scalar and vector structures needed for the mass and kinetic terms
that we will derive in Lecture 13.

\newpage
\section*{Questions for review}
\addcontentsline{toc}{section}{Questions for review}
\begin{enumerate}
\item Derive both Weyl generator matrices from \(J=L+R\) and
\(K=-i(L-R)\). Compare their action under a \(z\) rotation and boost.
\item Reproduce the trace proof excluding nontrivial finite-dimensional
unitary Lorentz representations. Why does it not exclude unitary particle states?
\item Why can parity act on the direct sum of the Weyl representations
but not within one of them alone?
\item Prove \(A^T E A=(\det A)E\), and derive the dual and conjugate-dual laws.
\item Show that the fundamental is self-dual but inequivalent to its
complex conjugate. Which part of the argument depends on boosts?
\item Compute \(B(\xi,\chi)\) with explicit raised indices.
Why does raising an undotted index not change its chirality?
\item Decompose a general two-index spinor into symmetric and
antisymmetric parts. What changes for \(\xi_\alpha\xi_\beta\)?
\item Verify that \(zz^\dagger\) gives a future null vector for a nonzero
commuting spinor. Why is a general Hermitian matrix not one such product?
\item Decompose a rank-two Lorentz tensor into irreducibles and check
the dimensions. Explain the reality relation between its two-form pieces.
\item Derive the sigma anticommutator and trace identities from the
Pauli multiplication rule. Check the vector bilinear for a \(z\) boost.
\item Compare the invariance of \(\xi^T E\chi\), \(\xi^\dagger\chi\),
and \(r^\dagger\xi\). State each field's transformation law.
\item Derive the exchange and Hermitian-conjugation rules for Grassmann
bilinears. Which rules change for ordinary complex wavefunctions?
\item Solve the massless Weyl equation at fixed positive-energy momentum.
Which spin polarization survives in the stated convention?
\end{enumerate}

\section*{Sources and reading}
\addcontentsline{toc}{section}{Sources and reading}
This lecture consolidates the Weyl-representation and contraction
material of the supplied \emph{Lecture 10: Representations of the Lorentz
Group} (2025), pp.\ 21--23, with \emph{Lecture 11: Spinors of the Lorentz
Group}, pp.\ 6--14, and the notation review in \emph{Lecture 12:
Action for Spinors}, pp.\ 1--7. The action derivations continue in
course Lecture 13. A useful independent reference is H.~K.~Dreiner,
H.~E.~Haber, and S.~P.~Martin,
\href{https://arxiv.org/abs/0812.1594}{arXiv:0812.1594}.
Translate epsilon, boost, and chirality conventions before comparing formulas.
\end{document}
